\section{Detecção: do erro ao alarme}
\label{sec:deteccao}

Temos um medidor de anomalia (o erro de reconstrução). Falta a regra que
transforma um número contínuo em uma decisão ``anômalo ou não''. Essa regra é o
\textbf{limiar} (\emph{threshold}).

\subsection{Erro de reconstrução por janela}
Para cada janela de 30 passos, comparamos a reconstrução com a entrada e
resumimos a diferença em um número --- usamos o \textbf{MAE} (Seção~\ref{sec:nn}),
robusto a \emph{outliers}:

\begin{lstlisting}[caption={Erro de reconstrução (MAE) por janela.}]
import numpy as np
recon = modelo.predict(X)          # X: (n_janelas, 30, 1)
erro  = np.mean(np.abs(recon - X), axis=(1, 2))   # um valor por janela
\end{lstlisting}

\subsection{Limiar estático: uma linha fixa}
\begin{definicao}[Limiar estático]
Um valor fixo, calculado como o \textbf{percentil 95 do erro no treino}.
Qualquer janela de teste com erro acima dele é marcada como anômala.
\end{definicao}

Calculá-lo \emph{no treino} (e não no teste) evita vazamento: o limiar reflete a
normalidade aprendida, não o período sob avaliação.

\begin{intuicao}
É como definir ``febre'' por um termômetro travado em 37,8\,°C. Simples e
estável --- mas não se adapta. Se a pessoa naturalmente vive mais quente por um
período (mudança de regime), o termômetro fixo apita o tempo todo.
\end{intuicao}

\subsection{Limiar dinâmico: uma linha que se adapta}
\begin{definicao}[Limiar dinâmico]
O percentil 95 recalculado continuamente sobre uma \textbf{janela móvel} dos
erros recentes (252 dias úteis $\approx$ 1 ano). Acompanha o regime local de
volatilidade.
\end{definicao}

O ponto delicado é a \textbf{causalidade}: a janela móvel só pode usar o
\emph{passado}. Se incluísse o próprio ponto (ou o futuro), uma anomalia ajudaria
a definir o limiar que deveria detectá-la --- vazamento de novo. Garantimos a
causalidade com um deslocamento de uma posição (\texttt{shift(1)}): cada ponto é
avaliado contra a estatística dos dias \emph{anteriores} a ele.

\begin{lstlisting}[caption={Limiar dinâmico causal (só passado), com pandas.}]
import pandas as pd
s = pd.Series(erro)
# shift(1) exclui o proprio ponto -> janela usa apenas o passado
limiar_din = s.shift(1).rolling(252, min_periods=252).quantile(0.95)
\end{lstlisting}

\begin{figure}[h]
\centering
\includegraphics[width=\textwidth]{m4_erro_limiares}
\caption{Erro de reconstrução no teste (2020--2024) com o limiar estático
(tracejado vermelho) e o dinâmico (verde). Onde o regime muda muito (AMER3), o
estático ``satura'' e marca quase tudo; o dinâmico se recalibra. Em ativos
estáveis os dois quase coincidem.}
\label{fig:teoria-m4}
\end{figure}

\begin{naprat}
Comparamos os dois limiares (uma das contribuições do projeto). No AMER3, após a
fraude, o erro fica permanentemente alto: o estático marca $\approx$27\% das
janelas (perde utilidade), enquanto o dinâmico cai para $\approx$14\% ao tratar o
novo patamar como o ``normal recente''. Em PETR4/VALE3 (estáveis) ambos ficam em
$\approx$3\%. Estudamos ainda o tamanho da janela móvel: 126 dias fica
ruidoso, 504 super-suaviza (marca quase nada), e \textbf{252} é o equilíbrio.
\end{naprat}
